\subsection{算子的积分}

设 G 与 H 为 R 上的两个 Hausdorff 局部凸空间，且现在假定 F 是由 G 到 H 中的连续线性映射所成的空间 \(\mathcal{L}(\mathrm{G};\mathrm{H})\) ，赋以逐点收敛拓扑。F 上的每个连续线性形式都可以延拓为乘积空间 \(H^{G}\) 上的连续线性形式（TVS, II, §4, No. 1, 命题 2），因而可以写成 \(u \mapsto \sum_{i=1}^{n} \langle u(\mathbf{a}_{i}), \mathbf{b}_{i}^{\prime} \rangle\) ，其中诸 \(a_{i}\)（相应地，诸 \(b_{i}^{\prime}\) ）是 G 的元素（相应地，H 的对偶 \(H^{\prime}\) 的元素）。说 T 到 F 中的映射 U 是标量本质 \(\mu\)-可积的，意思是对每个 \(a \in G\) 与每个 \(b \in H^{\prime}\) ，数值函数 \(t \mapsto \langle U(t) \cdot a, b^{\prime} \rangle\) 都是本质 \(\mu\)-可积的。

命题 9. --- 设 U 是标量本质 \(\mu\)-可积的、T 到 \(\mathbf{F} = \mathcal{L}_{s}(\mathbf{G}; \mathbf{H})\) 中的映射。为使 \(\int U d\mu \in F\) ，必须且只需下列两个条件成立：

\begin{enumerate}
\def\labelenumi{\alph{enumi})}
\item
  对每个 \(\mathbf{x} \in \mathbf{G}\) ， \(\int (U(t) \cdot \mathbf{x}) d\mu(t) \in \mathrm{H}\) 。
\item
  设 \(B'\) 为 \(H'\) 的任意等度连续子集；由线性形式 \(u_{y'} : x \mapsto \int \langle U(t) \cdot x, y' \rangle d\mu(t)\) （其中 \(y'\) 跑遍 \(B'\) ）所成的集是等度连续的。
\end{enumerate}

条件 a) 与 b) 是必要的。因为，对每个 \(\mathbf{x} \in G\) ，映射 \(\widetilde{\mathbf{x}}: V \mapsto V \cdot \mathbf{x}\) （由 \(\mathcal{L}_s(\mathrm{G};\mathrm{H})\) 到 H 中）是线性的，所以可以看出（No. 1，

命题 1）：\(\widetilde{\mathbf{x}}\circ U:t\mapsto U(t)\cdot \mathbf{x}\) 是标量本质 \(\mu\)-可积的，并且

\[
S \cdot \mathbf {x} = \int (U (t) \cdot \mathbf {x}) d \mu (t),\tag{1}
\]

其中 \(S = \int U d\mu \in \mathcal{L}_s(\mathrm{G};\mathrm{H})\) 。这就证明了 a) 。此外， (1) 式也可以写成

\[
\langle S \cdot \mathbf {x}, \mathbf {y} ^ {\prime} \rangle = \int \left\langle U (t) \cdot \mathbf {x}, \mathbf {y} ^ {\prime} \right\rangle d \mu (t) = \left\langle \mathbf {x}, u _ {\mathbf {y} ^ {\prime}} \right\rangle ,\tag{2}
\]

换言之 \(^{t}S \cdot y' = u_{y'}\) 。由于 S 连续， \(^{t}S\) 把 \(H'\) 的每个等度连续子集变换成 \(G'\) 的等度连续子集，由此得 b) 。

反之，设 a) 与 b) 成立。根据 a) ，公式 (1) 定义了一个由 G 到 H 中的线性映射 S ，且对每个 \(y' \in H'\) ，该映射满足 (2) （No. 1, 命题 1）；而这时条件 b) 表明 S 连续（TVS, II, §6, No. 4, 命题 5 与命题 6, 以及 III, §3, No. 5, 命题 7），因此 \(S \in \mathcal{L}_{s}(G;H)\) 。最后，公式 (2) 证明 \(S = \int U d\mu\) 。

推论. --- 命题 9 的条件 b) 在下列两种情形的每一种中都成立：

\(1^{\circ}\) 测度 \(\mu\) 有界，且若 S 为其支集，则 \(U(S)\) 是 \(\mathcal{L}(G;H)\) 的等度连续子集。

\(2^{\circ}\) 命题 9 的条件 a) 成立，空间 G 是桶式的，且对 T 的每个紧子集 K ， \(U(\mathrm{K})\) 是 \(\mathcal{L}_s(\mathrm{G};\mathrm{H})\) 的有界子集。

先考察情形 \(1^{\circ}\) 。我们可以限于 S = T 的情形（Ch. V, §7, No. 1, 定理 1）。于是，设 \(B'\) 为 \(H'\) 的任意等度连续子集，存在一个等度连续、凸、平衡且弱闭的子集 \(A' \subset G'\) ，使得 \({}^{t}U(t) \cdot y' \in A'\) 对一切 \(y' \in B'\) 与一切 \(t \in T\) 成立（TVS, II, §6, No. 4, 命题 6）。由于 U 是标量 \(\mu\)-可积的，映射 \(t \mapsto {}^{t}U(t) \cdot y'\) （取值于 G 的对偶 \(G'\) ，后者赋有 \(\sigma(G', G)\) ）是标量 \(\mu\)-可积的，于是可以写出

\[
u _ {\mathbf {y} ^ {\prime}} = \int \left(^ {t} U (t) \cdot \mathbf {y} ^ {\prime}\right) d \mu (t).
\]

由于 \(A'\) 对 \(\sigma(G', G)\) 是凸的且紧的，第 2 目命题 5 的推论表明 \(u_{\mathbf{y}'} \in \mu(T)A'\) （对每个 \(\mathbf{y}' \in B'\) ），这就证明了我们的断言。

现在考察情形 \(2^{\circ}\) 。对每个 \(y' \in H'\) 与 T 的每个紧子集 K ，令

\[
u _ {\mathrm{K}, \mathbf {y} ^ {\prime}} = \int \varphi_ {\mathrm{K}} (t) \left(^ {t} U (t) \cdot \mathbf {y} ^ {\prime}\right) d \mu (t),
\]

它是代数对偶 \(\mathbf{G}^*\) （ \(\mathbf{G}\) 的代数对偶）中的元素。由于 \(\mathbf{G}\) 是桶式的， \(\mathcal{L}_s(\mathbf{G};\mathbf{H})\) 的每个有界子集都等度连续（TVS, III, §4, No. 2, 定理 1）；把论证的第一部分应用于函数 \(\varphi_{\mathbf{K}}U\) 和有界测度 \(\varphi_{\mathbf{K}}\cdot \mu\) ，便得 \(u_{\mathbf{K},\mathbf{y}'}\in \mathbf{G}'\) 。此外，对拓扑 \(\sigma (\mathbf{G}^{*},\mathbf{G})\) ，有 \(u_{\mathbf{y}'} = \lim_{\mathbf{K}} = u_{\mathbf{K},\mathbf{y}'}\) ，极限是沿 \(\mathbf{T}\) 的紧子集所成的递增有向集取的（Ch. V, §1, No. 3, 命题 10）。为验证命题 9 的条件 b) ，根据命题 9，只需证明线性映射 \(S\) （由 (1) 定义，把 \(\mathbf{G}\) 映到 \(\mathbf{H}\) 中）是连续的；再者， \(\mathbf{G}\) 是桶式的，故只需证明 \(S\) 在 \(\mathbf{G}\) 与 \(\mathbf{H}\) 赋以它们的弱化拓扑时是连续的（TVS, IV, §1, No. 3, 命题 7）；最后，根据 (2) ，归结为证明： \(u_{\mathbf{y}'}\in \mathbf{G}'\) 对每个 \(\mathbf{y}'\in \mathbf{H}'\) 成立 。由于 \(u_{\mathbf{y}'}\) ——对拓扑 \(\sigma (\mathbf{G}^{*},\mathbf{G})\) 而言——属于集 \(\mathbf{M}'\) （由诸 \(u_{\mathbf{K},\mathbf{y}'}\) 组成，其中 \(\mathbf{K}\) 跑遍 \(\mathbf{T}\) 的紧子集所成之集）的闭包，只需证明 \(\mathbf{M}'\) 等度连续（TVS, III, §3, No. 4, 命题 4）；又由于 \(\mathbf{G}\) 是桶式的，这等价于说：对每个 \(\mathbf{x}\in \mathbf{G}\) ，诸 \(\langle \mathbf{x},u_{\mathbf{K},\mathbf{y}'}\rangle\) 所成之集有界（TVS, III, §4, No. 2, 定理 1）。而这由关系式

\[
| \langle \mathbf {x}, u _ {\mathrm{K}, \mathbf {y} ^ {\prime}} \rangle | = \left| \int \varphi_ {\mathrm{K}} (t) \langle U (t) \cdot \mathbf {x}, \mathbf {y} ^ {\prime} \rangle d \mu (t) \right| \leqslant \int | \langle U (t) \cdot \mathbf {x}, \mathbf {y} ^ {\prime} \rangle | d \mu (t).
\]

命题 10. --- 设 U 是 T 到 F = \(\mathcal{L}_s(\mathrm{G};\mathrm{H})\) 中的映射。在下列三种情形的每一种中， U 都是标量本质 \(\mu\)-可积的，且 \(\int U d\mu \in \mathcal{L}_s(\mathrm{G};\mathrm{H})\) ：

\begin{enumerate}
\def\labelenumi{\alph{enumi})}
\item
  H 拟完备， \(\mu\) 有界，且若 S 为其支集，则 U 是 \(\mu\)-可测的且 \(U(S)\) 等度连续。
\item
  H 半自反， \(\mu\) 有界，且若 S 为其支集，则 U 是标量 \(\mu\)-可测的且 \(U(S)\) 等度连续。
\item
  H 半自反， G 桶式， U 标量本质 \(\mu\)-可积，且对 T 的每个紧子集 K ， \(U(\mathrm{K})\) 有界。
\end{enumerate}

在三种情形中， U 标量本质可积这一事实都是明显的；根据命题 9 及其推论，在每种情形只需验证命题 9 的条件 a) 。而在第一种情形，这一条件由第 2 目命题 8 得出，在后两种情形则由第 2 目命题 7 的推论得出。

